\section{Магнетизм}
\subsection{Магнитное поле}

%%%%%%%%%%%%%%%%%%%%%%%%%%%%%%%%%%%%%%%%%%%%%%%%%%%%
%%%%%%%%%%%%%%%%%%%%%%%%%%%%%%%%%%%%%%%%%%%%%%%%%%%%
 \z{Найти индукцию магнитного поля в центре кольца радиуса $R$, по которому протекает ток $I$.}
%
{$B = \frac{\mu_0 I}{2R}$}

%%%%%%%%%%%%%%%%%%%%%%%%%%%%%%%%%%%%%%%%%%%%%%%%%%%%
 \z{Длинные прямые провода с током пересекаются под прямым углом. Определите индукцию магнитного поля на биссектрисе прямого угла.}
%
{$0$ или $\frac{\mu_0 (I_1 + I_2)}{2\pi h}$}

%%%%%%%%%%%%%%%%%%%%%%%%%%%%%%%%%%%%%%%%%%%%%%%%%%%%
 \z{Во сколько раз уменьшится индукция магнитного поля в центре кольца с током, если его согнуть под углом $90^\circ$? Ток в кольце не меняется.}
%
{$\frac{\sqrt{2}}{2}$}

%%%%%%%%%%%%%%%%%%%%%%%%%%%%%%%%%%%%%%%%%%%%%%%%%%%%
 \z[{\begin{wrapfigure}[6]{r}{0.15\textwidth}%
		\vspace{-15pt}%
		\begin{flushright}%
 			\includegraphics[width=0.15\textwidth]{5.1.4}%
		\end{flushright}%
	\end{wrapfigure}%
	}]%
{По проводнику, расположенному в одной плоскости, как изображено на рисунке, течёт ток. Найти индукцию магнитного поля в произвольной точке линии $AB$.}
%
{$0$}

%%%%%%%%%%%%%%%%%%%%%%%%%%%%%%%%%%%%%%%%%%%%%%%%%%%%
 \z{Сила тока в проводнике, согнутом под прямым углом, $I = 15\,\text{А}$. Какой будет индукция магнитного поля в точке, лежащей на биссектрисе того угла и отстоящей от вершины на расстоянии $r = 0.05\,\text{м}$?}
%
{$B = \frac{\mu_0 I(1+\sqrt{2})}{2\pi r}$}



%%%%%%%%%%%%%%%%%%%%%%%%%%%%%%%%%%%%%%%%%%%%%%%%%%%%
 \z[{\begin{wrapfigure}[5]{r}{0.30\textwidth}%
		\vspace{-15pt}%
		\begin{flushright}%
 			\includegraphics[width=0.30\textwidth]{5.1.6}%
		\end{flushright}%
	\end{wrapfigure}%
	}]%
{Вычислить индукцию магнитного поля в точке $O$, если проводник имеет вид, показанный на рисунке (а). Сила тока $I$ в проводнике и радиус $R$ известны.}
%
{$B = \frac{\mu_0 I}{4R}$}

%%%%%%%%%%%%%%%%%%%%%%%%%%%%%%%%%%%%%%%%%%%%%%%%%%%%
 \z{То же, но для проводника, изображённого на рисунке (б).}
%
{$B = \frac{3\mu_0 I}{8R} + \frac{\mu_0 I}{2\pi R}$}







